%! Solving Rational Equations — Unit 4, Lesson 4.6 — Student Packet Cover Sheet
\documentclass[10pt]{article}
\usepackage{algebra2-article}
\usepackage{algebra2-boxes}
\usepackage{ltablex}
\keepXColumns

\begin{document}

% Full-bleed forest banner
\begin{tikzpicture}[remember picture, overlay]
  \fill[forest] (current page.north west)
    rectangle ([yshift=-0.9in]current page.north east);
\end{tikzpicture}
\vspace{-0.8in}

\begin{center}
  {\color{white}\textbf{\LARGE Algebra 2}}\\[4pt]
  {\color{forestmid}\large\textbf{Unit 4 \quad Rational Functions}}\\[6pt]
  {\color{white}\textbf{\large Lesson 4.6 \quad Solving Rational Equations}}
\end{center}

\vspace{0.22in}
\namedateperiod
\vspace{0.18in}

\begin{learningtargetbox}
\small
\begin{enumerate}[leftmargin=1.5em, itemsep=5pt]
  \item \ldots{} \textbf{solve} a rational equation by factoring the denominators, writing the
        \textbf{restriction list first}, multiplying \emph{every term} by the LCD, and solving the
        linear or quadratic equation that is left.
  \item \ldots{} \textbf{check} every candidate against that list, discard the ones that match, and
        \textbf{verify} the survivors --- algebraically and by reading a graph.
  \item \ldots{} \textbf{explain why} a candidate can be extraneous: multiplying both sides by an
        expression that could be zero is a step that cannot be undone.
  \item \ldots{} \textbf{write} a rational equation to model a real situation, solve it, and say what
        the answer means --- including which answers to reject, and whether the reason is algebraic or
        contextual.
\end{enumerate}
\end{learningtargetbox}

\vspace{0.14in}

\begin{tocbox}
\small
\renewcommand{\arraystretch}{1.5}
\begin{tabularx}{\linewidth}{@{} c l X r @{}}
\rowcolor{forestbg}
\textbf{\#} & \textbf{Component} & \textbf{Description} & \textbf{Score} \\
\midrule
1 & Warm-Up        & Spiral review: restrictions, the LCD, and what happens when you multiply an equation by zero & \blank{1.2cm} \\
2 & Guided Notes   & Expression vs.\ equation, the four steps, why extraneous solutions happen, two worked examples, and cross-multiplying & \blank{1.2cm} \\
3 & Group Activity & \emph{Solve, Then Check} --- twin equations, a graphical check, an error analysis, a work-rate model, and a design task, by tier & \blank{1.2cm} \\
4 & Exit Ticket    & Quick check: one full solve, one written justification, plus an SOL-style item & \blank{1.2cm} \\
5 & Homework       & Restriction lists, six solves, candidate checks, a graph to read, an error analysis, and a river-current model & \blank{1.2cm} \\
\midrule
  & \multicolumn{2}{r}{\textbf{Total}} & \blank{1.2cm} \\
\end{tabularx}
\end{tocbox}

\vspace{0.10in}

\begin{remindbox}
\small All unit long you have found excluded values and kept them \emph{off} a graph. Today they come
back as answers --- the ones you have to throw away. \textbf{Four steps, always in this order.}
\textbf{(1) Factor and restrict:} factor every denominator and write the list of excluded values
\emph{at the top of your work}, before you touch the equation. \textbf{(2) Multiply by the LCD:} every
single term, including any term that has no denominator at all --- you are multiplying both
\emph{sides}, not just the fractions. \textbf{(3) Solve} the linear or quadratic equation that is left,
with the toolkit you already own. \textbf{(4) Check:} compare each candidate to your Step 1 list. Any
candidate that appears on the list is \textbf{extraneous} --- discard it. Verify the survivors in the
\emph{original} equation. \textbf{Why any of this is necessary:} multiplying both sides by a nonzero
number can always be undone by dividing, so it never changes the answers. But the LCD is an
\emph{expression}, and at each excluded value that expression equals zero --- and multiplying by zero
cannot be undone. That is the whole mechanism, and it has a consequence worth memorizing: an extraneous
solution is never a random number. It is always a value from your Step 1 list. On a graph, it is even
easier to see. Move everything to one side, combine, and the true solutions are the
\textbf{$x$-intercepts} while the extraneous candidates are the \textbf{holes} --- Lesson 4.5's cancel
test, doing today's job. And remember that ``no solution'' is a complete answer, not a sign you did
something wrong.
\end{remindbox}

\end{document}
